
\begin{frame}[t]{From one score to a sum of rules}
\lead{Linear and logistic regression start from one score: $F(x)=\beta_0+x^\top\beta$.}
\begin{columns}[T,onlytextwidth]
\column{.47\textwidth}
\begin{definitionbox}\textbf{Regression}\par
Prediction: $\widehat y=F(x)$.
\[\ell(y,F)=\tfrac12(y-F)^2.\]
\end{definitionbox}
\column{.49\textwidth}
\begin{definitionbox}\textbf{Binary classification}\par
Prediction: $\widehat p=\sigm(F(x))$, $y\in\{0,1\}$.
\[\ell(y,F)=\log(1+e^F)-yF.\]
\end{definitionbox}
\end{columns}
\vspace{.15cm}
\uncover<2->{\begin{takeaway}
Learn the score as a sum of fitted rules:
\[F_T(x)=F_0(x)+\sum_{t=1}^T\alpha_t h_t(x).\]
Each rule can be simple; their combination can express a richer function.
\end{takeaway}}
\note{The sigmoid is s(z)=1/(1+exp(-z)). An intercept is included in the linear score. Feature expansion is another way to enrich a linear predictor; ensemble learning instead fits and combines multiple component functions. An ensemble is the umbrella term; boosting and bagging are two different construction strategies. This slide previews gradient boosting, where the previous regression or classification loss can be retained. AdaBoost first gives a particularly clean derivation using exponential loss and minus/plus-one labels; explicitly announce that change. Sums of linear functions of exactly the same features remain linear, so expressive gains require suitable nonlinear or varied rules, for example threshold functions. A weak model here means a simple component with limited individual predictive ability; the formal weighted-edge assumption will be stated when proving the AdaBoost training bound.}
\end{frame}

\begin{frame}[t]{Learning Goals}
By the end of the lecture, you should be able to:
\begin{itemize}\setlength{\itemsep}{.24cm}
\item explain the high-level idea of ensemble learning
\item explain how AdaBoost combines simple rules into a stronger classifier;
\item connect residual fitting to gradient descent in prediction space;
\end{itemize}
\end{frame}
